\documentclass{standalone}
\usepackage{circuitikz}
\usetikzlibrary{calc}
\definecolor{circuitnotemark}{HTML}{FE00FE}
\begin{document}
\begin{circuitikz}[american, line width=0.8pt]
\ctikzset{bipoles/length=1.2cm}
\foreach \x in {0,3,6,9,12} {\foreach \y in {0,-3,-6} {\fill[gray, opacity=0.35] (\x,\y) circle (1.65pt);}}
\node[gray, font=\scriptsize] at (0,1.26) {1};
\node[gray, font=\scriptsize] at (3,1.26) {2};
\node[gray, font=\scriptsize] at (6,1.26) {3};
\node[gray, font=\scriptsize] at (9,1.26) {4};
\node[gray, font=\scriptsize] at (12,1.26) {5};
\node[gray, font=\scriptsize, anchor=east] at (-1.26,0) {a};
\node[gray, font=\scriptsize, anchor=east] at (-1.26,-3) {b};
\node[gray, font=\scriptsize, anchor=east] at (-1.26,-6) {c};
\coordinate (a1) at (0,0);
\coordinate (a5) at (12,0);
\coordinate (b1) at (0,-3);
\coordinate (b5) at (12,-3);
\coordinate (bp5-1) at (0,-4.5);
\coordinate (bp5-5) at (12,-4.5);
\coordinate (c1) at (0,-6);
\coordinate (c5) at (12,-6);
\draw (a1) node[ocirc]{} node[above left]{IN}; % line 3
\draw (a1) to[R, l_=$R_{1}$, a^=$1\,\mathrm{k}\Omega$] (a5); % line 4
\draw (b1) to[R, l_=$R_{2}$, a^=$2\,\mathrm{k}\Omega$] (b5); % line 5
\draw (bp5-1) to[R, l_=$R_{3}$, a^=$3\,\mathrm{k}\Omega$] (bp5-5); % line 6
\draw (c1) to[R, l_=$R_{4}$, a^=$4\,\mathrm{k}\Omega$] (c5); % line 7
\draw (a5) node[ocirc]{} node[above left]{OUT}; % line 8
\draw (a1) -- (b1); % line 10
\draw (b1) -- (c1); % line 10
\draw (a5) -- (b5); % line 11
\draw (b5) -- (c5); % line 11
\node[circ] at (a1) {};
\node[circ] at (a5) {};
\node[circ] at (b1) {};
\node[circ] at (b5) {};
\node[circ] at (bp5-1) {};
\node[circ] at (bp5-5) {};
\path (15,-8.869) rectangle (23.302,0.593); % line 13
\node[anchor=west, inner sep=0, circuitnotemark, font=\large] at (15,0) {X}; % line 13
\node[anchor=west, inner sep=0, circuitnotemark, font=\large] at (15,-0.487) {X}; % line 13
\node[anchor=west, inner sep=0, circuitnotemark, font=\large] at (15,-0.974) {X}; % line 13
\node[anchor=west, inner sep=0, circuitnotemark, font=\large] at (15,-1.46) {X}; % line 13
\node[anchor=west, inner sep=0, circuitnotemark, font=\large] at (15,-1.947) {X}; % line 13
\node[anchor=west, inner sep=0, circuitnotemark, font=\large] at (15,-2.434) {X}; % line 13
\node[anchor=west, inner sep=0, circuitnotemark, font=\large] at (15,-2.921) {X}; % line 13
\node[anchor=west, inner sep=0, circuitnotemark, font=\large] at (15,-3.408) {X}; % line 13
\node[anchor=west, inner sep=0, circuitnotemark, font=\large] at (15,-3.895) {X}; % line 13
\node[anchor=west, inner sep=0, circuitnotemark, font=\large] at (15,-4.381) {X}; % line 13
\node[anchor=west, inner sep=0, circuitnotemark, font=\large] at (15,-4.868) {X}; % line 13
\node[anchor=west, inner sep=0, circuitnotemark, font=\large] at (15,-5.355) {X}; % line 13
\node[anchor=west, inner sep=0, circuitnotemark, font=\large] at (15,-5.842) {X}; % line 13
\node[anchor=west, inner sep=0, circuitnotemark, font=\large] at (15,-6.329) {X}; % line 13
\node[anchor=west, inner sep=0, circuitnotemark, font=\large] at (15,-6.816) {X}; % line 13
\node[anchor=west, inner sep=0, circuitnotemark, font=\large] at (15,-7.302) {X}; % line 13
\node[anchor=west, inner sep=0, circuitnotemark, font=\large] at (15,-7.789) {X}; % line 13
\node[anchor=west, inner sep=0, circuitnotemark, font=\large] at (15,-8.276) {X}; % line 13
\node[anchor=south west, inner sep=2pt, circuitnotemark, font=\LARGE\bfseries] (circuittitle) at (current bounding box.north west) {X};
\path (circuittitle.south west) rectangle ++(10.516,0.853);
\end{circuitikz}
\end{document}
